\documentclass[11pt]{article}
\usepackage[T1]{fontenc}
\usepackage[utf8]{inputenc}
\usepackage{lmodern}
\showboxdepth=10000 \showboxbreadth=1000000 \errorcontextlines=-1
\makeatletter
\newwrite\mtdump
\def\mtdumpfont{\edef\mtname{\expandafter\string\the\font}%
  \immediate\write\mtdump{@@font \mtname\space quad=\number\fontdimen6\font\space size=\number\pdffontsize\font}%
  \count@=0 \loop
  \immediate\write\mtdump{\number\count@\space\number\fontcharwd\font\count@\space\number\lpcode\font\count@\space\number\rpcode\font\count@\space\number\efcode\font\count@}%
  \advance\count@ 1 \ifnum\count@<256 \repeat}
\def\mtprobe#1{%
  \setbox4=\vbox{\pdfprotrudechars=0 \hsize=500pt \parindent=0pt \parfillskip=0pt \leftskip=0pt \rightskip=0pt \pretolerance=-1 \tolerance=10000 #1M\par}\showbox4
  \setbox4=\vbox{\pdfprotrudechars=0 \hsize=1pt \parindent=0pt \parfillskip=0pt \leftskip=0pt \rightskip=0pt \pretolerance=-1 \tolerance=10000 #1M\par}\showbox4
  {#1\mtdumpfont}}
\makeatother
\begin{document}
\immediate\openout\mtdump=\jobname.dump
\setbox2=\vbox{\hsize=\maxdimen \pretolerance=-1 \parindent=0pt %
Self-evident, state-of-the-art, and/or long-term: these
compound words, together with dis\-cre\-tion\-ary hyphens, test how pdfTeX
treats hyphens at the right margin. Yes---really. Why? Because ``T'' and
``V'' and ``W'' and ``Y'' all protrude, as do 1, 4 and 7.\par}\showbox2
\setbox0=\vbox{\hsize=80pt \parindent=0pt \emergencystretch=30pt \hbadness=500 \tracingparagraphs=1 \tracingonline=0 %
\immediate\write\mtdump{@@params hsize=\number\hsize\space pretolerance=\number\pretolerance\space tolerance=\number\tolerance\space emergencystretch=\number\emergencystretch\space linepenalty=\number\linepenalty\space hyphenpenalty=\number\hyphenpenalty\space exhyphenpenalty=\number\exhyphenpenalty\space adjdemerits=\number\adjdemerits\space doublehyphendemerits=\number\doublehyphendemerits\space finalhyphendemerits=\number\finalhyphendemerits\space looseness=\number\looseness\space lastlinefit=\number\lastlinefit\space parindent=\number\parindent\space leftskip=\the\leftskip\space rightskip=\the\rightskip\space parfillskip=\the\parfillskip\space pdfadjustspacing=\number\pdfadjustspacing\space pdfprotrudechars=\number\pdfprotrudechars\space hbadness=\number\hbadness\space hfuzz=\number\hfuzz}%
Self-evident, state-of-the-art, and/or long-term: these
compound words, together with dis\-cre\-tion\-ary hyphens, test how pdfTeX
treats hyphens at the right margin. Yes---really. Why? Because ``T'' and
``V'' and ``W'' and ``Y'' all protrude, as do 1, 4 and 7.
\showlists\par}\showbox0
\mtprobe{\normalfont}\mtprobe{\itshape}\mtprobe{\bfseries}
\immediate\closeout\mtdump
\end{document}
